\section{Removal squares and repair children}
\label{sec:area_law_simultaneous_repair}

The removal step in \cite[Section~11]{OpenAI2026AreaLaw} takes place
before replacement regions are created. Every old region loses the same
union of closed squares, including squares whose interiors will be left
unassigned at the stopping scale.

% Source: 10-geometry.tex, geometry:trim-radius, lines 380--399;
% geometry:birth-separation, lines 584--591.
% Isolated-star induction and replacement regions are not asserted here.

\begin{definition}[Dyadic removal square]
    \label{def:area_law_repair_square}
    \lean{TNLean.PEPS.AreaLaw.Geometry.repairHalfSide,
        TNLean.PEPS.AreaLaw.Geometry.repairSquare}
    \leanok
    \uses{thm:area_law_exponent_gaps}
    Retain the fixed exponent $\rho_0=(1+\beta)/2$. For a positive
    parent scale $S$ and a mark $v\in\mathbb R^2$, define
    \begin{align}
        d(S)   & =2^{\lfloor\rho_0\log_2 S\rfloor}, &
        Q(v,S) & =v+[-d(S),d(S)]^2.
        \label{eq:area_law_repair_square}
    \end{align}
    The floor and the exponent are integers. Equivalently, $Q(v,S)$
    is the closed ball of radius $d(S)$ for the maximum metric.
\end{definition}

% Source: 10-geometry.tex, geometry:trim-margins, lines 386--391.
\begin{theorem}[Removal margins]
    \label{thm:area_law_repair_margins}
    \lean{TNLean.PEPS.AreaLaw.Geometry.rho_pos_lt_one,
        TNLean.PEPS.AreaLaw.Geometry.repairHalfSide_le_rpow,
        TNLean.PEPS.AreaLaw.Geometry.half_rpow_le_repairHalfSide,
        TNLean.PEPS.AreaLaw.Geometry.repairHalfSide_lt_self,
        TNLean.PEPS.AreaLaw.Geometry.eventually_repairHalfSide_le,
        TNLean.PEPS.AreaLaw.Geometry.eventually_repairHalfSide_margins}
    \leanok
    \uses{def:area_law_repair_square}
    The exponent satisfies $0<\rho_0<1$. For every $S>0$,
    \begin{equation}
        \tfrac12S^{\rho_0}\le d(S)\le S^{\rho_0},
        \label{eq:area_law_repair_margins}
    \end{equation}
    and $d(S)<S$ whenever $S>1$. For every $a_0>0$ there is $S_0$ such
    that every $S\ge S_0$ satisfies $d(S)\le a_0S/1000$. Hence a stopping
    scale $S_*$ with the margins of \cite[Section~11]{OpenAI2026AreaLaw}
    exists and may be taken as large as required.
\end{theorem}

\begin{proof}
    \leanok
    \uses{thm:area_law_exponent_gaps}
    Put $x=\rho_0\log_2 S$, so that $S^{\rho_0}=2^x$ and
    $d(S)=2^{\lfloor x\rfloor}$. The inequalities
    $\lfloor x\rfloor\le x<\lfloor x\rfloor+1$ give
    \eqref{eq:area_law_repair_margins}. The exponent gaps give
    $\rho_0<1$, so $S^{\rho_0}<S$ for $S>1$. Finally
    $S^{\rho_0}=S^{\rho_0-1}S$ and $S^{\rho_0-1}\to0$ as $S\to\infty$.
\end{proof}

\begin{definition}[Simultaneous old-region remnants]
    \label{def:area_law_simultaneous_remnants}
    \lean{TNLean.PEPS.AreaLaw.Geometry.repairRemovalSet,
        TNLean.PEPS.AreaLaw.Geometry.simultaneousRepairRemnants}
    \leanok
    \uses{def:area_law_repair_square}
    Let $(v_j,S_j)_{j\in J}$ be a finite family of marks, each with a
    location $v_j\in\mathbb R^2$ and a positive assigned scale $S_j$;
    distinct indices may share a location. Let $(R_i)_{i\in I}$ be
    the current actual regions of the old identifiers. Set
    \begin{align}
        Q_J  & =\bigcup_{j\in J}Q(v_j,S_j), &
        R'_i & =R_i\setminus Q_J.
        \label{eq:area_law_simultaneous_remnants}
    \end{align}
    This formula applies to every $i\in I$. The labels and remembered
    birth regions retain their values at the removal step.
\end{definition}

% Source: OpenAI2026AreaLaw, Section 11, 10-geometry.tex,
% lines 401--409 and geometry:repair-induction, lines 506--512.
% This definition records child indices, locations and scales only.
% It asserts no injectivity, isolated-star invariant or repair induction.
\begin{definition}[Children of nonterminal parents]
    \label{def:area_law_repair_children}
    \lean{TNLean.PEPS.AreaLaw.Geometry.repairChildOffset,
        TNLean.PEPS.AreaLaw.Geometry.RepairChildIndex,
        TNLean.PEPS.AreaLaw.Geometry.repairChildLocation,
        TNLean.PEPS.AreaLaw.Geometry.repairChildScale}
    \leanok
    \uses{def:area_law_repair_square}
    Let $(v_i,S_i)_{i\in I}$ be the current parent family, with
    $v_i\in\mathbb R^2$ and $S_i>0$, and retain the stopping threshold
    $S_*$. Set $I_+=\{i\in I:d(S_i)\ge S_*\}$ and
    $P=\{-1,0,1\}^2$. The next children are indexed by $I_+\times P$.
    For each $(i,p)\in I_+\times P$, their locations and assigned scales are
    \begin{align}
        v'_{i,p} & =v_i+d(S_i)p, & S'_{i,p} & =d(S_i).
        \label{eq:area_law_repair_children}
    \end{align}
    These offsets give the center, corners and side midpoints of
    $Q(v_i,S_i)$. Each child retains its parent and offset as indices;
    coincident locations do not identify these indices. A parent with
    $d(S_i)<S_*$ contributes no children, and its square is terminal.
    This is the marking rule of \cite[Section~11]{OpenAI2026AreaLaw}.
\end{definition}

\subsection{Depth and descendant count of a repair branch}
\label{sec:area_law_repair_depth}

Every child of a parent of scale $S$ receives the scale $d(S)$. The marks
of generation $j$ descended from one initial mark of scale $S$ therefore
all have the scale $d^{j}(S)$, the $j$-fold iterate of $d$, and each
generation has at most $9^j$ marks. The following estimates concern these
iterated scales only; they do not construct the replacements.

% Source: OpenAI2026AreaLaw, Section 11, 10-geometry.tex, lines 668--686;
% OpenAI2026PolynomialPEPS, Lemma 3.2, 02-information.tex, lines 302--311.
\begin{theorem}[Depth of a power-contracting recursion]
    \label{thm:area_law_power_recursion_depth}
    \lean{TNLean.PEPS.AreaLaw.Geometry.log_iterate_le_of_le_rpow,
        TNLean.PEPS.AreaLaw.Geometry.one_le_log_add_exp_one,
        TNLean.PEPS.AreaLaw.Geometry.exists_iterate_depth_bound,
        TNLean.PEPS.AreaLaw.Geometry.exists_iterate_generation_count_bound,
        TNLean.PEPS.AreaLaw.Geometry.mul_rpow_le_max_mul_rpow,
        TNLean.PEPS.AreaLaw.Geometry.forall_le_iterate_of_le_iterate}
    \leanok
    Let $0<\rho<1$ and $S_*>1$, and let $f$ satisfy $0<f(u)\le u^{\rho}$
    for every $u\ge S_*$. A recursion from a scale $S>0$ passes from
    generation $i$ to generation $i+1$ only when $f^{i}(S)\ge S_*$. If the
    generations $0,\dots,m-1$ have scale at least $S_*$, then
    $\log f^{i}(S)\le\rho^{i}\log S$ for $i\le m$, and there is a constant
    $C>0$, independent of $S$ and $m$, with
    \begin{align}
        m & \le C+C\log\log(S+e).
        \notag
    \end{align}
    For every $b\ge0$ there is moreover $C>0$ with
    $\sum_{j=0}^{m}b^{j}\le C\bigl(\log(S+e)\bigr)^{C}$ under the same
    condition. If $0<f(u)\le u^{\rho}$ for every $u>0$ and $S_*\ge1$, a
    generation of scale at least $S_*$ has all earlier generations of scale
    at least $S_*$.
\end{theorem}

\begin{proof}
    \leanok
    Induction on $i$ gives the logarithmic bound. If $m\ge1$, then
    $0<\log S_*\le\rho^{m-1}\log S$, so $\log S>0$ and
    $(m-1)\log(1/\rho)\le\log\log S-\log\log S_*$; moreover
    $\log\log S\le\log\log(S+e)$, which is nonnegative. For the count, with
    $B=b+2$ the sum is at most $B^{m+1}$, and the depth bound gives
    $B^{m+1}\le B^{C+1}\bigl(\log(S+e)\bigr)^{C\log B}$. Enlarging the
    exponent only increases a power of a number at least one. For the last
    assertion, a scale $x<S_*$ has $f(x)\le x^{\rho}<S_*^{\rho}\le S_*$, so
    once a generation falls below $S_*$ all later ones do.
\end{proof}

% Source: OpenAI2026AreaLaw, Section 11, 10-geometry.tex, lines 668--675.
\begin{theorem}[Depth of a repair branch]
    \label{thm:area_law_repair_depth}
    \lean{TNLean.PEPS.AreaLaw.Geometry.repairHalfSide_pos,
        TNLean.PEPS.AreaLaw.Geometry.log_repairHalfSide_le,
        TNLean.PEPS.AreaLaw.Geometry.log_repairHalfSide_iterate_le,
        TNLean.PEPS.AreaLaw.Geometry.exists_repair_depth_bound}
    \leanok
    \uses{def:area_law_repair_square}
    For every $S>0$ and $j\ge0$,
    \begin{align}
        \log d^{j}(S) & \le\rho_0^{j}\log S.
        \label{eq:area_law_repair_depth_log}
    \end{align}
    For every $S_*>1$ there is a constant $C>0$ such that, for every $S>0$
    and every $j\ge0$ with $d^{j}(S)\ge S_*$,
    \begin{align}
        j & \le C+C\log\log\bigl(e^{e}+S\bigr).
    \end{align}
\end{theorem}

\begin{proof}
    \leanok
    \uses{thm:area_law_repair_margins, thm:area_law_power_recursion_depth}
    The half-side $d(S)$ is a power of two and $d(S)\le S^{\rho_0}$, so
    $\log d(S)\le\rho_0\log S$. Induction on $j$ gives
    (\ref{eq:area_law_repair_depth_log}). Since $0<d(u)\le u^{\rho_0}$ for
    every $u>0$, Theorem~\ref{thm:area_law_power_recursion_depth} shows that
    if $d^{j}(S)\ge S_*$, then all generations $0,\dots,j$ have scale at least
    $S_*$, and it bounds $j+1$ by $C+C\log\log(S+e)$. Finally
    $\log\log(S+e)\le\log\log(e^{e}+S)$.
\end{proof}

% Source: OpenAI2026AreaLaw, Section 11, 10-geometry.tex, lines 676--686,
% geometry:descendant-count.
\begin{theorem}[Descendant count of an initial mark]
    \label{thm:area_law_repair_descendant_count}
    \lean{TNLean.PEPS.AreaLaw.Geometry.exists_repair_descendant_count_bound,
        TNLean.PEPS.AreaLaw.Geometry.exists_repair_descendant_count_layer_bound}
    \leanok
    \uses{def:area_law_repair_square}
    For every $S_*>1$ there is a constant $C>0$ such that, for every $S>0$
    and every $m\ge0$ with $d^{m}(S)\ge S_*$,
    \begin{align}
        \sum_{j=0}^{m}9^{j} & \le C\bigl(\log(S+2)\bigr)^{C}.
    \end{align}
    For every $M>0$ there is moreover a constant $C>0$, depending only on
    $S_*$ and $M$, such that the same sum is at most $C(k+1)^{C}$ whenever
    $S\le M2^{k}$ for a nonnegative integer $k$.
\end{theorem}

\begin{proof}
    \leanok
    \uses{thm:area_law_power_recursion_depth}
    If $d^{m}(S)\ge S_*$, all generations $0,\dots,m$ have scale at least
    $S_*$, by Theorem~\ref{thm:area_law_power_recursion_depth}; in particular
    $S\ge S_*>1$. The same theorem with $b=9$ bounds the sum by
    $C\bigl(\log(S+e)\bigr)^{C}$. Since $S+e\le e(S+2)$ and
    $\log(S+2)\ge1$, we have $\log(S+e)\le2\log(S+2)$, and enlarging the
    constant gives the first assertion. For the second assertion,
    $S+2\le(M+2)2^{k}$ gives $\log(S+2)\le(\log(M+2)+1)(k+1)$.
\end{proof}

% Source: OpenAI2026AreaLaw, Section 11, 10-geometry.tex, lines 761--768,
% geometry:fine-clearance.
\begin{theorem}[Fine clearance threshold]
    \label{thm:area_law_fine_clearance}
    \lean{TNLean.PEPS.AreaLaw.Geometry.eventually_fine_clearance}
    \leanok
    Let $c_1>0$ and $C_B\ge0$. Every sufficiently large real number $u$
    satisfies
    \begin{align}
        c_1u^{\rho_0} & >2\lceil C_Bu\rceil^{\beta}.
    \end{align}
\end{theorem}

\begin{proof}
    \leanok
    \uses{thm:area_law_exponent_gaps}
    For $u\ge1$, $\lceil C_Bu\rceil\le(C_B+1)u$, so the right side is at
    most $2(C_B+1)^{\beta}u^{\beta}$. The exponent gaps give $\beta<\rho_0$,
    and $u^{\beta-\rho_0}$ tends to zero.
\end{proof}
